% Generated draft: requires mathematical and source review.
\documentclass[11pt]{article}
\usepackage[margin=25mm]{geometry}
\usepackage{fontspec}
\setmainfont{Latin Modern Roman}
\usepackage{amsmath,amssymb,mathtools,unicode-math}
\setmathfont{Latin Modern Math}
\usepackage{enumitem,needspace,xurl,hyperref}
\setlength{\parindent}{0pt}
\setlength{\parskip}{5pt}
\setlength{\emergencystretch}{3em}
\newcommand{\N}{\mathbb{N}}
\newcommand{\Z}{\mathbb{Z}}
\newcommand{\Q}{\mathbb{Q}}
\newcommand{\R}{\mathbb{R}}
\newcommand{\C}{\mathbb{C}}
\newcommand{\F}{\mathbb{F}}
\newcommand{\A}{\mathbb{A}}
\newcommand{\PP}{\mathbb{P}}
\newcommand{\E}{\mathbb{E}}
\newcommand{\eps}{\varepsilon}
\newcommand{\dd}{\,\mathrm{d}}
\begin{document}
\section*{UnsolvedMath ID: 30004729}
\textbf{Source problem code:} catalogue entry
\subsection{Definitions and mathematical statement}
The supplied record concerns a numerical scheme for the general Monge--Amp\`ere equation in the far-field refractor setting. Here, ``numerical solutions'' means solutions of that scheme. The record supplies neither the equation nor the scheme, so these terms identify the source's intended objects without specifying their mathematical definitions.

The source question is: Do the numerical solutions of this scheme converge?

This is an unresolved formulation rather than a fully specified convergence assertion. The record does not give the optical geometry, domains, coefficients, prescribed data, boundary or other admissibility conditions, discrete equations, refinement parameter, or selection of discrete solutions. It also does not specify a convergence topology, a limiting object, or whether convergence is required for every family of discrete solutions or only along suitable subsequences. Consequently, a precise quantified convergence statement cannot be recovered from the supplied evidence without additional source information. In particular, no specific Monge--Amp\`ere equation or convergence theorem is asserted here.
\subsection{Short English statement}
The question asks whether the numerical solutions of the intended far-field refractor scheme converge. The supplied record omits the scheme, hypotheses, limiting target, and meaning of convergence, so a precise convergence assertion requires further source information. Its open status is supplied curation, not independently verified.
\subsection{Sources}
\begin{itemize}
\item Supplied local UnsolvedMath registry record, permanent ID \(30004729\), title ``Convergence of Far-Field Refractor Schemes,'' source problem code \texttt{catalogue entry}, in the Oberwolfach Reports Open Problems collection. No named proposer is supplied; the recorded proposed year is 2021.
\item Supplied report citation: \emph{Numerical Methods for Fully Nonlinear and Related PDEs} (2022), Oberwolfach Reports, DOI 10.4171/owr/2021/31, pp.\ 1651--1674. \url{https://doi.org/10.4171/owr/2021/31}. The record describes Oberwolfach Reports as documenting mathematical workshops and their open-problem discussions. The report itself has not been independently consulted for this draft.
\item Supplied literature source: ``Oberwolfach Reports article 7155452.'' \url{https://ems.press/journals/owr/articles/7155452}.
\item Literature assessment dated exactly \texttt{2026-08-22}: the supplied curation labels the problem open and states that convergence remains open, with no general convergence theorem resolving the unspecified scheme. This is supplied dated curation derived from the OWR record, not an independently verified literature assessment; time-sensitive status and cited claims require independent verification.
\item The supplied statement status is \texttt{corrected\_verified}. Its accompanying verification note says that only clear formula delimiters, notation, grammar, and typography were repaired, preserving mathematical content and scope. This metadata does not supply the missing formulation or establish mathematical correctness.
\end{itemize}
\end{document}
